\question{1.4}{Geef voor de volgende variabelen aan of deze een nominale, ordinale, interval- of ratioschaal heeft:

    \begin{enumerate}[label=(\alph*)]
        \item de industrietak waarin werknemers een baan hebben
        \item de speelduur van een DVD
        \item de kleur van tulpen
        \item de jaaromzet (in euro) van bedrijven
        \item het aantal sterren dat de moeilijkheidsgraad van puzzelboekjes aangeeft
        \item de hoogte boven de zeespiegel van wintersportdorpen
    \end{enumerate}
}

\answer{
    \begin{enumerate}[label=(\alph*)]
        \item Industrietak: nominale variabele
        \item Speelduur DVD: ratiovariabele
        \item Kleur tulpen: nominale variabele
        \item Jaaromzet: ratiovariabele
        \item Aantal sterren moeilijkheidsgraad: ordinale variabele
        \item Hoogte boven zeespiegel: ratiovariabele
    \end{enumerate} 
}